\begin{textsample}{POS Dim 2 – to – Score 21.00 – to/to\_em\_69.txt}  \label{ex:f2_pos_140}
11.
ARQUITETURAS PARA O ENSINO MÉDIO E AS POSSIBILIDADES DE ESTRUTURA CURRICULAR A Lei nº 13.415/2017 alterou a Lei de \textbf{Diretrizes} e Bases da Educação Nacional e estabeleceu mudanças na estrutura do Ensino Médio, ampliando o tempo mínimo do estudante na escola, de 800 horas para 1.000 horas anuais ( até 2022 ) e definindo uma nova organização curricular, mais \textbf{flexível}, que contemple uma Base Nacional Comum Curricular ( BNCC ) e a \textbf{oferta} de diferentes possibilidades de escolhas aos estudantes : os \textbf{Itinerários} \textbf{Formativos} ( no \textbf{regime} parcial ) e, nas Escolas de Tempo Integral, as Unidades Curriculares Integradoras com foco nas áreas de conhecimento e na formação \textbf{técnica} e \textbf{profissional}.
A mudança tem como objetivos garantir a \textbf{oferta} de educação de \textbf{qualidade} a todos os jovens brasileiros e de aproximar as escolas à realidade dos estudantes do século XXI, considerando as novas demandas e complexidades do mundo do trabalho e da vida em sociedade.
O infográfico a seguir traz possíveis arranjos com distribuição da \textbf{carga} \textbf{horária} da Formação Geral Básica e dos \textbf{Itinerários} \textbf{Formativos}, ao longo dos três anos do Ensino Médio.
Ponto de atenção : É importante destacar que estas possibilidades de Arquitetura limitam -se ao período de 2022 a 2024, considerando que a \textbf{legislação} \textbf{vigente} estabelece que, a partir de 2024, a \textbf{carga} \textbf{horária} anual passará a ser de 1.400 horas ( relógio ).
Para a \textbf{rede} \textbf{estadual}, orienta -se que seja \textbf{destinada} uma \textbf{carga} \textbf{horária} específica para o desenvolvimento do Projeto de Vida dos estudantes desde o início do Ensino Médio, para que estes tenham a oportunidade de exercer seu protagonismo, considerando que esta é a etapa em que ocorre o maior número de evasão. 11.1 Possibilidades de \textbf{oferta} de \textbf{carga} \textbf{horária} para o Ensino Médio A arquitetura curricular do território do Tocantins compreenderá, de 2022 a 2024, a seguinte \textbf{carga} \textbf{horária} : - Formação Geral Básica com, no máximo, 1.800 horas ; e - \textbf{Itinerários} \textbf{Formativos} com, no mínimo, 1.200 horas.
A \textbf{Portaria} nº 521, de 13 de julho de 2021, \textbf{instituiu} o Cronograma Nacional de Implementação do Novo Ensino Médio, tendo como objetivo apoiar as Unidades da Federação no processo de implementação dos documentos curriculares, \textbf{alinhados} à Base Nacional Comum Curricular ( BNCC, 2018 ), bem como efetivar a operacionalização do art. 24, § 1º, e do art. 36 da Lei nº 9.394, de 20 de dezembro de 1996.
Desta forma, de 2022 a 2024, as Escolas de Ensino Médio deverão \textbf{ofertar} \textbf{carga} \textbf{horária} de mil horas anuais.
A implementação do novo referencial curricular no Tocantins \textbf{atenderá} ao seguinte cronograma, conforme Portaria/MEC nº 521/2021 : - 2022 — Implementação do Documento Curricular na 1ª série do Ensino Médio. - 2023 — Implementação do Documento Curricular nas 1ª e 2ª séries do Ensino Médio. - 2024 — Implementação do Documento Curricular nas 3ª séries do Ensino Médio.
O cronograma contempla também o monitoramento da implementação do Documento Curricular do Território do Tocantins e formação continuada dos \textbf{profissionais} da educação, nos anos de 2022 a 2024.

% matched lemmas: alinhar, atender, carga, destinar, diretriz, estadual, flexível, formativo, horário, instituir, itinerário, legislação, oferta, ofertar, portaria, profissional, qualidade, rede, regime, técnico, vigente
\end{textsample}
